\chapter{Supplementary Material}
\label{app:supplementary}

\section{NPP indicator definitions}
\label{app:indicators}

\Cref{tab:indicators} lists the fourteen Node--Place--People indicators used in the
prioritization weighting (\cref{sec:w4}). The vitality proxy is excluded from the
weighting for the reason given in \cref{sec:w4}. Each indicator is stored both raw and
normalized to $[0,1]$; the normalization rule depends on the variable's distribution.
Right-skewed count and economic variables are transformed with $\log(1+x)$ before
min--max scaling to prevent a few large values from compressing the rest; bounded ratios
are min--max scaled directly. The marginalization index is additionally
\emph{inverted} ($1-\text{minmax}$) so that higher values denote greater need, and its
missing \ageb{}s are median-imputed rather than zero-filled (\cref{sec:area-integrity}).

\begin{table}[htbp]
\centering
\caption{The fourteen NPP indicators, by dimension, with normalization rule.}
\label{tab:indicators}
\begin{tabularx}{\linewidth}{llX}
\toprule
Dimension & Indicator & Normalization \\
\midrule
Node   & \code{n\_intersections}         & $\log$+minmax \\
Node   & \code{n\_intersection\_density} & $\log$+minmax \\
Node   & \code{n\_street\_density}       & $\log$+minmax \\
\midrule
Place  & \code{p\_poi\_density}          & $\log$+minmax \\
Place  & \code{p\_employment\_proxy}     & $\log$+minmax \\
Place  & \code{p\_retail\_density}       & $\log$+minmax \\
Place  & \code{p\_service\_density}      & $\log$+minmax \\
Place  & \code{p\_land\_use\_mix}        & minmax \\
\midrule
People & \code{pe\_population}           & $\log$+minmax \\
People & \code{pe\_pop\_density}         & $\log$+minmax \\
People & \code{pe\_marginacion}          & minmax, inverted \\
People & \code{pe\_rezago}               & minmax \\
People & \code{pe\_dep\_ratio}           & minmax \\
People & \code{pe\_youth\_share}         & minmax \\
\bottomrule
\end{tabularx}
\end{table}

\section{Reproducibility}
\label{app:reproducibility}

The analysis is organized as a sequence of workstream orchestrators (W1 through W8) that
run against a PostgreSQL/PostGIS database, with a from-scratch bootstrap that builds the
schema, loads the committed inputs, and materializes the feature table before any
workstream runs. All computation uses the canonical projected reference system
EPSG:6372; all credentials and constants are centralized so that no value is duplicated
across scripts. The transfer cities are driven by the same modules through a
city-selection flag, reusing the pure \zmg functions unchanged.

Two reproducibility notes bear on the numbers in this thesis. First, the feature table is
re-derivable from committed census, economic, and network inputs; a from-scratch rebuild
reproduces the place, people, and equity columns almost exactly, while the three
network-derived indicators drift slightly because the street graph is pulled live and can
change between builds. Second, all headline figures in this thesis were computed against
the corrected \num{1881}-\ageb universe and the corrected equity term; any earlier figures
for the same quantities, computed before the integrity corrections of
\cref{sec:area-integrity}, are superseded.
